\documentclass[12pt, a4paper]{report}
\usepackage[utf-8]{inputenc}
\usepackage[vietnamese]{babel}
\usepackage{xcolor}
\usepackage{listings}
\usepackage{geometry}
\usepackage{fancyhdr}
\usepackage{hyperref}
\usepackage{graphicx}
\usepackage{tcolorbox}
\usepackage{array}
\usepackage{booktabs}

\geometry{margin=2.5cm}

\lstset{
    basicstyle=\ttfamily\small,
    breaklines=true,
    columns=fullflexible,
    backgroundcolor=\color{gray!10},
    frame=single
}

\pagestyle{fancy}
\fancyhf{}
\rhead{\thepage}
\lhead{Hệ Thống Quản Lý Khách Sạn}

\title{\textbf{GIẢI THÍCH CHI TIẾT VỀ CHỨC NĂNG, TASK VÀ USER STORY}\\
\textbf{Hệ Thống Quản Lý Khách Sạn Toàn Diện}\\
\large (Hotel Management System)}

\author{AI Task Generator}
\date{\today}

\begin{document}

\maketitle

\tableofcontents
\newpage

\chapter{GIỚI THIỆU TỔNG QUAN}

\section{Về Hệ Thống Quản Lý Khách Sạn}

Hệ thống quản lý khách sạn là một ứng dụng web-based toàn diện, giúp quản lý các hoạt động hàng ngày của nhà hàng. Hệ thống bao gồm các chức năng chính như:

\begin{itemize}
    \item Quản lý phòng và tính khả dụng
    \item Quản lý đặt phòng (booking)
    \item Quản lý khách hàng
    \item Quản lý thanh toán và hóa đơn
    \item Quản lý nhân viên
    \item Báo cáo và phân tích dữ liệu
    \item Tích hợp với các hệ thống bên ngoài
\end{itemize}

\section{Thông Số Kỹ Thuật Chính}

\begin{itemize}
    \item \textbf{Nền tảng:} Web-based với Responsive Design
    \item \textbf{Backend:} Java/Spring Boot
    \item \textbf{Database:} PostgreSQL
    \item \textbf{Deployment:} Docker, Kubernetes
    \item \textbf{Server:} AWS hoặc Azure Cloud
    \item \textbf{Thời gian phát triển:} 8 tháng
    \item \textbf{Ngân sách:} \$150,000
    \item \textbf{Nhân lực:} 1 PM + 3 Developer
\end{itemize}

\newpage

\chapter{PHÂN TÍCH CÁC CHỨC NĂNG CHÍNH}

\section{1. Quản Lý Phòng (Room Management)}

\subsection{Mô Tả Chức Năng}

Chức năng này cho phép quản lý các thông tin chi tiết về phòng trong khách sạn, bao gồm loại phòng, trạng thái, tiện nghi, và giá cả.

\subsection{Các Task Liên Quan}

\begin{tcolorbox}[colback=blue!5, colframe=blue!50]
\textbf{Task 1: Quản lý các loại phòng với giá và tiện nghi khác nhau}
\begin{itemize}
    \item \textbf{Mức độ ưu tiên:} High
    \item \textbf{Story Points:} 5 SP
    \item \textbf{Danh mục:} Hotel Management
    \item \textbf{Phạm vi:} Happy Path, Edge Cases, CRUD Operations
\end{itemize}
\end{tcolorbox}

\subsection{User Story Ví Dụ}

\begin{tcolorbox}[colback=green!5, colframe=green!50]
\textbf{User Story: Quản lý Phòng và Tài sản - Happy Path}

\begin{lstlisting}
As a Manager
I want to quản lý phòng và tài sản 
So that I can maintain hotel assets and room information

Acceptance Criteria:
✓ Given: Đã đăng nhập với vai trò Quản lý
  When: Thực hiện quản lý phòng và tài sản
  Then: Hệ thống xử lý thành công 
        và hiển thị kết quả

✓ Given: Dữ liệu đầu vào không hợp lệ
  When: Cố gắng thực hiện hành động
  Then: Hệ thống hiển thị thông báo lỗi rõ ràng
\end{lstlisting}
\end{tcolorbox}

\subsection{Chi Tiết Các Acceptance Criteria}

\begin{table}[h]
\centering
\begin{tabular}{|p{3cm}|p{4cm}|p{3cm}|}
\hline
\textbf{Given (Điều kiện)} & \textbf{When (Khi)} & \textbf{Then (Thì)} \\
\hline
Đã đăng nhập với vai trò Quản lý & Thực hiện quản lý phòng & Hệ thống lưu thông tin \\
\hline
Dữ liệu không hợp lệ & Gửi form không hợp lệ & Hiển thị lỗi rõ ràng \\
\hline
Phòng đã tồn tại & Thêm phòng mới & Hiển thị thông báo trùng \\
\hline
\end{tabular}
\end{table}

\section{2. Quản Lý Đặt Phòng (Booking Management)}

\subsection{Mô Tả Chức Năng}

Cho phép khách hàng đặt phòng, kiểm tra tính khả dụng, xác nhận, hủy hoặc sửa đổi đơn đặt phòng.

\subsection{Các Task Chính}

\begin{tcolorbox}[colback=orange!5, colframe=orange!50]
\textbf{Task 2.1: Phải cho phép đặt phòng mới}
\begin{itemize}
    \item Story Points: 5 SP
    \item Scope: Happy Path + Edge Cases
    \item Input Required: check-in date, check-out date, room type, guest info
\end{itemize}

\textbf{Task 2.2: Phải kiểm tra tính khả dụng của phòng}
\begin{itemize}
    \item Story Points: 5 SP
    \item Scope: Happy Path + Edge Cases
    \item Check: theo loại phòng và ngày yêu cầu
\end{itemize}

\textbf{Task 2.3: Phải cho phép xác nhận, hủy và chỉnh sửa đơn}
\begin{itemize}
    \item Story Points: 5 SP
    \item Scope: CRUD Operations
    \item Access: Khách hàng và Quản lý
\end{itemize}
\end{tcolorbox}

\subsection{Quy Trình Booking - Sơ Đồ Luồng}

\begin{tcolorbox}[colback=gray!10]
\begin{verbatim}
KHÁCH HÀNG
   |
   v
[1. Tìm kiếm phòng trống] 
   |
   +---> Input: check-in, check-out, room type
   |
   v
[2. Hệ thống hiển thị phòng khả dụng + giá]
   |
   v
[3. Khách hàng cung cấp thông tin cá nhân]
   |
   +---> Name, Email, Phone, ID
   |
   v
[4. Tạo đơn đặt phòng trong hệ thống]
   |
   v
[5. Gửi email xác nhận cho khách]
   |
   v
[QUẢN LÝ xác nhận và xử lý thanh toán]
\end{verbatim}
\end{tcolorbox}

\section{3. Quản Lý Thanh Toán (Payment Management)}

\subsection{Các Componet Chính}

\begin{table}[h]
\centering
\begin{tabular}{|l|p{5cm}|p{4cm}|}
\hline
\textbf{Chức Năng} & \textbf{Mô Tả} & \textbf{Story Points} \\
\hline
Tạo hóa đơn & Tạo hóa đơn chi tiết cho mỗi lần lưu trú & 5 SP \\
\hline
Phương thức thanh toán & Hỗ trợ tiền mặt, thẻ, chuyển khoản & 5 SP \\
\hline
Thanh toán một phần & Cho phép đặt cọc hoặc trả góp & 5 SP \\
\hline
Lịch sử thanh toán & Lưu trữ và truy vấn lịch sử giao dịch & 5 SP \\
\hline
Xuất hóa đơn điện tử & Export PDF, Email, Print & 5 SP \\
\hline
\end{tabular}
\end{table}

\newpage

\chapter{CÁCH VIẾT USER STORY HIỆU QUẢ}

\section{1. Định Nghĩa User Story}

\begin{tcolorbox}[colback=purple!5, colframe=purple!50]
\textbf{User Story} là một mô tả ngắn gọn về một tính năng từ góc độ người dùng cuối.

Cấu trúc cơ bản:
\begin{lstlisting}
As a [User Type]
I want to [Action/Feature]
So that [Business Value/Benefit]

Acceptance Criteria:
• Given [Precondition]
  When [Action]
  Then [Expected Result]
\end{lstlisting}
\end{tcolorbox}

\section{2. Các Loại User (Personas)}

\subsection{Trong Hệ Thống Quản Lý Khách Sạn}

\begin{enumerate}
    \item \textbf{Khách hàng (Customer)}
    \begin{itemize}
        \item Tìm kiếm phòng
        \item Đặt phòng
        \item Thanh toán
        \item Xem lịch sử lưu trú
    \end{itemize}
    
    \item \textbf{Lễ tân (Receptionist)}
    \begin{itemize}
        \item Xử lý check-in/check-out
        \item Quản lý phòng
        \item Hỗ trợ khách hàng
    \end{itemize}
    
    \item \textbf{Quản lý (Manager)}
    \begin{itemize}
        \item Quản lý toàn hệ thống
        \item Xem báo cáo
        \item Phân quyền người dùng
        \item Tính toán doanh thu
    \end{itemize}
    
    \item \textbf{Nhân viên vệ sinh (Housekeeping)}
    \begin{itemize}
        \item Xem lịch dọn phòng
        \item Báo cáo tình trạng phòng
    \end{itemize}
    
    \item \textbf{Admin (Administrator)}
    \begin{itemize}
        \item Quản lý hệ thống
        \item Cấu hình tham số
        \item Bảo mật
    \end{itemize}
\end{enumerate}

\section{3. Cấu Trúc Acceptance Criteria (AC)}

\subsection{Mô Hình Gherkin - Phổ Biến Nhất}

\begin{tcolorbox}[colback=cyan!5, colframe=cyan!50]
\begin{lstlisting}
Scenario: [Tên kịch bản rõ ràng]

Given [Điều kiện ban đầu]
  And [Điều kiện thêm]
  
When [Hành động/Sự kiện xảy ra]
  And [Hành động phụ]
  
Then [Kết quả mong đợi]
  And [Kết quả thêm]
\end{lstlisting}
\end{tcolorbox}

\subsection{Ví Dụ Cụ Thể - Booking}

\begin{tcolorbox}[colback=yellow!5, colframe=yellow!50]
\begin{lstlisting}
Scenario: Khách hàng đặt phòng thành công

Given: Khách hàng đã truy cập vào trang booking
  And: Có các phòng trống cho ngày cần đặt
  
When: Khách hàng chọn loại phòng
  And: Nhập ngày check-in và check-out
  And: Cung cấp thông tin cá nhân đầy đủ
  And: Nhấn nút "Xác nhận đặt phòng"
  
Then: Hệ thống tạo booking record mới
  And: Gửi email xác nhận tới khách hàng
  And: Cập nhật trạng thái phòng thành "Đã đặt"
  And: Hiển thị thông báo "Đặt phòng thành công"
\end{lstlisting}
\end{tcolorbox}

\newpage

\section{4. Phân Loại User Story Theo Kích Thước}

\begin{table}[h]
\centering
\begin{tabular}{|l|p{3cm}|p{3cm}|p{3cm}|}
\hline
\textbf{Kích Thước} & \textbf{Story Points} & \textbf{Thời Gian} & \textbf{Ví Dụ} \\
\hline
Rất nhỏ & 1-2 SP & 2-4 giờ & Thêm field mới \\
\hline
Nhỏ & 3-5 SP & 1-2 ngày & CRUD đơn giản \\
\hline
Vừa & 8-13 SP & 1 tuần & Check-in/out \\
\hline
Lớn & 20-40 SP & 2-3 tuần & Booking system \\
\hline
Rất lớn & >40 SP & >1 tháng & Tích hợp OTA \\
\hline
\end{tabular}
\end{table}

\section{5. Tiêu Chí User Story Tốt (INVEST)}

\begin{tcolorbox}[colback=red!5, colframe=red!50]
\textbf{I} - Independent: Độc lập, không phụ thuộc vào story khác

\textbf{N} - Negotiable: Có thể thảo luận chi tiết

\textbf{V} - Valuable: Mang lại giá trị cho người dùng

\textbf{E} - Estimable: Có thể ước tính kích thước

\textbf{S} - Small: Đủ nhỏ để hoàn thành trong 1 sprint

\textbf{T} - Testable: Có thể kiểm thử được
\end{tcolorbox}

\newpage

\chapter{CẤU TRÚC TASK VÀ CÁC CẤP ĐỘ}

\section{1. Phân Loại Task Theo Cấp Độ}

\subsection{1.1 Task Happy Path}

\begin{tcolorbox}[colback=green!5]
\textbf{Happy Path:} Quy trình hoạt động bình thường, cách tử dùng hành động chính xác

\textbf{Ví dụ:} 
\begin{lstlisting}
Task: Quản lý Đặt phòng - Happy Path

Scenario: Khách hàng đặt phòng và thanh toán 
thành công

Steps:
1. Khách hàng đăng nhập thành công
2. Tìm kiếm phòng khả dụng
3. Chọn phòng và ngày
4. Nhập thông tin khách
5. Thanh toán
6. Nhận xác nhận

Expected: Tất cả bước thực hiện thành công
\end{lstlisting}
\end{tcolorbox}

\subsection{1.2 Task Edge Cases & Validation}

\begin{tcolorbox}[colback=orange!5]
\textbf{Edge Cases:} Các tình huống ranh giới, dữ liệu lệch

\textbf{Ví dụ:}
\begin{lstlisting}
Task: Quản lý Đặt phòng - Edge Cases & Validation

Test Cases:
1. Check-out date < Check-in date
   Expected: Hiển thị lỗi "Ngày không hợp lệ"

2. Email không đúng định dạng
   Expected: Yêu cầu nhập lại

3. Số phòng đặt vượt quá số phòng trống
   Expected: Thông báo không đủ phòng

4. Thanh toán với số tiền âm
   Expected: Từ chối giao dịch

5. Dữ liệu null/empty
   Expected: Yêu cầu điền đầy đủ
\end{lstlisting}
\end{tcolorbox}

\subsection{1.3 Task CRUD Operations}

\begin{tcolorbox}[colback=blue!5]
\textbf{CRUD:} Create (Tạo), Read (Đọc), Update (Sửa), Delete (Xóa)

\textbf{Ví dụ - 5 Scenarios:}
\begin{lstlisting}
1. CREATE - Tạo mới phòng
   Input: room_type, room_number, price, status
   Output: Phòng được lưu vào database

2. READ/View - Xem danh sách phòng
   Input: filter parameters (type, availability)
   Output: Danh sách phòng phù hợp

3. UPDATE - Sửa thông tin phòng
   Input: room_id, new_price, new_status
   Output: Thông tin được cập nhật

4. DELETE - Xóa phòng
   Input: room_id, reason
   Output: Phòng bị xóa (soft delete)

5. VIEW DETAILS - Xem chi tiết
   Input: room_id
   Output: Tất cả thông tin phòng
\end{lstlisting}
\end{tcolorbox}

\section{2. Story Points Ước Tính}

\begin{table}[h]
\centering
\begin{tabular}{|l|c|p{4cm}|}
\hline
\textbf{Mức Độ Phức Tạp} & \textbf{Story Points} & \textbf{Đặc Điểm} \\
\hline
Rất đơn giản & 1-2 & UI thay đổi, label mới \\
\hline
Đơn giản & 3-5 & CRUD 1 entity, validation cơ bản \\
\hline
Trung bình & 8-13 & Multi-entity, business logic \\
\hline
Phức tạp & 20-40 & Tích hợp, performance, report \\
\hline
Rất phức tạp & >40 & Nhiều hệ thống, machine learning \\
\hline
\end{tabular}
\end{table}

\newpage

\chapter{PHÂN TÍCH CHI TIẾT - HOTEL MANAGEMENT}

\section{1. Yêu Cầu Chức Năng (Functional Requirements)}

\subsection{FR-1: Room Management Module}

\begin{tcolorbox}[colback=green!5]
\begin{lstlisting}
Yêu cầu: Quản lý thông tin phòng

Database Schema:
- Rooms Table
  ├─ room_id (PK)
  ├─ room_number (Unique)
  ├─ room_type (Single/Double/Suite)
  ├─ status (Available/Booked/Cleaning)
  ├─ price (Decimal)
  ├─ max_guests (Integer)
  ├─ amenities (JSON)
  └─ created_at, updated_at

User Story:
As a Manager
I want to manage room details
So that I can maintain accurate room status

AC:
✓ Can create new room
✓ Can view room list
✓ Can update room info
✓ Can mark room as unavailable
✓ System validates required fields
\end{lstlisting}
\end{tcolorbox}

\subsection{FR-2: Booking Management Module}

\begin{tcolorbox}[colback=blue!5]
\begin{lstlisting}
Yêu cầu: Quản lý đặt phòng

Database Schema:
- Bookings Table
  ├─ booking_id (PK)
  ├─ customer_id (FK)
  ├─ room_id (FK)
  ├─ check_in_date
  ├─ check_out_date
  ├─ number_of_guests
  ├─ status (Pending/Confirmed/Cancelled)
  ├─ payment_status (Unpaid/Partial/Paid)
  └─ created_at, updated_at

- Customers Table
  ├─ customer_id (PK)
  ├─ name
  ├─ email
  ├─ phone
  ├─ id_number
  ├─ address
  └─ loyalty_points

User Story A:
As a Customer
I want to book a room
So that I can reserve accommodation

User Story B:
As a Manager
I want to manage all bookings
So that I can ensure operational efficiency
\end{lstlisting}
\end{tcolorbox}

\subsection{FR-3: Payment & Billing Module}

\begin{tcolorbox}[colback=orange!5]
\begin{lstlisting}
Yêu cầu: Quản lý thanh toán

Database Schema:
- Invoices Table
  ├─ invoice_id (PK)
  ├─ booking_id (FK)
  ├─ room_charge
  ├─ service_charges
  ├─ tax
  ├─ discount
  ├─ total_amount
  ├─ payment_method
  ├─ payment_status
  └─ payment_date

- Transactions Table (Audit Trail)
  ├─ transaction_id
  ├─ invoice_id
  ├─ amount
  ├─ timestamp
  └─ reference_number

User Story:
As a Customer
I want to pay my bill
So that I can settle my account

Acceptance:
✓ Multiple payment methods
✓ Receipt generation
✓ Email invoice
✓ Partial payment allowed
✓ Tax calculation correct
\end{lstlisting}
\end{tcolorbox}

\newpage

\section{2. Yêu Cầu Phi Chức Năng (Non-Functional Requirements)}

\begin{table}[h]
\centering
\begin{tabular}{|p{2cm}|p{5cm}|p{3cm}|}
\hline
\textbf{Loại} & \textbf{Yêu Cầu} & \textbf{Chỉ Số} \\
\hline
Performance & Response time endpoints & <2 giây \\
\hline
Availability & System uptime & 99.9\% \\
\hline
Scalability & Concurrent users & 100+ users \\
\hline
Security & Data encryption & SSL/TLS \\
\hline
Compliance & Data protection & GDPR compliant \\
\hline
Usability & User training & <4 hours \\
\hline
Backup & Daily backup & Hourly snapshot \\
\hline
\end{tabular}
\end{table}

\chapter{QUY TRÌNH PHÁT TRIỂN PROJECT}

\section{Timeline - 8 Tháng}

\begin{tcolorbox}[colback=gray!5]
\begin{lstlisting}
Phase 1: Phân tích yêu cầu (2 tuần)
├─ Tập hợp requirement từ stakeholders
├─ Phân tích functional & non-functional
└─ Tạo specification document

Phase 2: Thiết kế hệ thống (4 tuần)
├─ Database design
├─ API design
├─ UI/UX mockups
└─ Architecture review

Phase 3: Phát triển Core Modules (12 tuần)
├─ Room Management
├─ Booking System
├─ Payment Integration
├─ User Management
└─ Basic Reporting

Phase 4: Phát triển Advanced Features (8 tuần)
├─ OTA Integration (Booking.com, Expedia)
├─ Electronic lock integration
├─ Mobile app support
├─ Advanced Analytics
└─ AI recommendations

Phase 5: Testing & QA (4 tuần)
├─ Unit testing
├─ Integration testing
├─ User acceptance testing
└─ Performance testing

Phase 6: Deployment & Documentation (4 tuần)
├─ Production setup
├─ Staff training
├─ Data migration
└─ Go-live support
\end{lstlisting}
\end{tcolorbox}

\chapter{BEST PRACTICES VÀ LƯU Ý}

\section{1. Viết User Story Tốt}

\begin{enumerate}
    \item \textbf{Cụ thể:} Tránh quá chung chung
    \begin{itemize}
        \item ❌ Sai: "As a user, I want the system to work well"
        \item ✓ Đúng: "As a Manager, I want to view daily booking report, so that I can plan staff"
    \end{itemize}
    
    \item \textbf{Đo lường được:} Có acceptance criteria rõ ràng
    \begin{itemize}
        \item ❌ Sai: "System should be fast"
        \item ✓ Đúng: "Search should complete within 2 seconds"
    \end{itemize}
    
    \item \textbf{Giá trị rõ ràng:} Lợi ích cụ thể cho user
    \begin{itemize}
        \item ❌ Sai: "Add API endpoint"
        \item ✓ Đúng: "So that I can integrate with PMS"
    \end{itemize}
    
    \item \textbf{Kích thước vừa phải:} Hoàn thành trong 1 sprint
    \begin{itemize}
        \item ❌ Sai: "Xây dựng toàn bộ hệ thống"
        \item ✓ Đúng: "Implement user login for Manager role"
    \end{itemize}
\end{enumerate}

\section{2. Tổ Chức Task Hiệu Quả}

\begin{tcolorbox}[colback=yellow!5]
\begin{lstlisting}
Epic
  └─ Feature (Đặc tính)
      └─ User Story (User story)
          └─ Task (Công việc cụ thể)
              └─ Subtask (Chi tiết công việc)

Ví dụ cho Booking Feature:

EPIC: Hotel Booking System
├─ FEATURE: Customer Booking
│  ├─ US: Search & View Available Rooms
│  │  ├─ Task: Create search API
│  │  ├─ Task: Build search UI
│  │  └─ Task: Write unit tests
│  │
│  ├─ US: Complete Booking
│  │  ├─ Task: Create booking record
│  │  ├─ Task: Send confirmation email
│  │  └─ Task: Update room status
│  │
│  └─ US: Cancel/Modify Booking
│     ├─ Task: Implement cancel logic
│     ├─ Task: Refund processing
│     └─ Task: Audit trail logging
│
└─ FEATURE: Admin Booking Management
   ├─ US: View all bookings
   ├─ US: Manually create booking
   └─ US: Generate booking reports
\end{lstlisting}
\end{tcolorbox}

\section{3. Definition of Done (DoD)}

\begin{tcolorbox}[colback=red!5]
Một User Story được coi là hoàn thành khi:

\begin{enumerate}
    \item ✓ Code được viết và reviewed
    \item ✓ Unit tests có coverage ≥ 80\%
    \item ✓ Tất cả acceptance criteria passed
    \item ✓ Integration tests passed
    \item ✓ No critical/high bugs
    \item ✓ Documentation updated
    \item ✓ Code merged to master
    \item ✓ QA verified on staging
\end{enumerate}
\end{tcolorbox}

\newpage

\chapter{VÍ DỤ THỰC VẬN HỌC}

\section{Ví Dụ 1: Complete Booking Flow}

\begin{tcolorbox}[colback=cyan!5, colframe=cyan!50]
\textbf{EPIC: Complete Hotel Booking System}

\textbf{FEATURE F1: Search \& Book Rooms}

\textbf{User Story US1.1}
\begin{lstlisting}
Title: Search available rooms by date

As a Customer
I want to search for available rooms 
by check-in and check-out dates
So that I can find suitable accommodation

Acceptance Criteria:
Given: Customer is on search page
When: Enter check-in date (2024-05-01)
  And: Enter check-out date (2024-05-05)
  And: Select room type (Double)
  And: Click search button
Then: System should display available rooms
  And: Show room price for selected dates
  And: Display room amenities
  And: Show availability status
  And: Load within 2 seconds

Edge Cases:
Scenario: Invalid dates
- Given: Check-out before check-in
  When: User clicks search
  Then: Show error "Invalid date range"

Scenario: No rooms available
- Given: All rooms booked
  When: User searches
  Then: Show "No rooms available"
  And: Suggest alternative dates

Task Breakdown:
1. Backend: Create room search API
   - Query rooms by availability
   - Calculate pricing
   - Return room details
   
2. Frontend: Build search form
   - Date picker component
   - Room type dropdown
   - Display results
   
3. Testing:
   - Unit test search logic
   - Integration test API
   - UI test form validation

Story Points: 5 SP
Time Estimate: 3-4 days
\end{lstlisting}
\end{tcolorbox}

\section{Ví Dụ 2: Payment Processing}

\begin{tcolorbox}[colback=magenta!5, colframe=magenta!50]
\textbf{User Story US2.1}
\begin{lstlisting}
Title: Process payment for booking

As a Customer
I want to pay for my booking
So that I can confirm my reservation

Acceptance Criteria:

Scenario: Successful payment
Given: Customer completed booking
  And: Payment amount is 5,000,000 VND
When: Select payment method (Credit Card)
  And: Enter card details correctly
  And: Click confirm payment
Then: Payment processing starts
  And: Transaction successful message shown
  And: Customer receives confirmation email
  And: Booking status changes to "Confirmed"

Scenario: Insufficient funds
Given: Payment amount is 5,000,000
When: Card balance is only 2,000,000
  And: Click confirm payment
Then: Show error "Insufficient funds"
  And: Suggest alternative payment methods

Scenario: Partial payment
Given: Customer can pay 2,000,000 (40%)
When: Select "Pay partial"
  And: Enter 2,000,000
  And: Confirm
Then: Record partial payment
  And: Remaining amount: 3,000,000
  And: Set reminder for remaining payment
  And: Keep booking status "Pending"

Scenario: Payment timeout
Given: Payment gateway timeout
When: Process continues >30 seconds
Then: Cancel request gracefully
  And: Show error message
  And: Allow user to retry

Task Breakdown:
1. Backend:
   - Integrate payment gateway (Stripe/PayU)
   - Create transaction record
   - Send confirmation
   
2. Frontend:
   - Payment form
   - Error handling
   - Success confirmation
   
3. Security:
   - PCI compliance
   - Encrypt card data
   - SSL/TLS

Story Points: 8 SP
Dependencies: Booking system
Acceptance Criteria: 5 scenarios
\end{lstlisting}
\end{tcolorbox}

\newpage

\chapter{CÔNG CỤ VÀ TEMPLATE}

\section{User Story Template}

\begin{lstlisting}
Title: [Clear, specific title]

User Story:
As a [User Type]
I want to [Action]
So that [Business Value]

Background:
[Prerequisites and context]

Acceptance Criteria:

Scenario 1: [Happy path]
Given [Precondition 1]
  And [Precondition 2]
When [User action]
  And [System action]
Then [Expected result]

Scenario 2: [Edge case]
Given [Different condition]
When [Different action]
Then [Different result]

Additional Notes:
- Dependencies: [Related stories]
- Nice to have: [Optional features]
- Technical notes: [Implementation hints]

Estimation:
- Story Points: [ ]
- Ideal Hours: [ ]
- Complexity: [ ] Low / [ ] Medium / [ ] High

Attachments:
- Wireframes: [Links/Images]
- Mockups: [Links/Images]
- Related docs: [Links]
\end{lstlisting}

\section{Task Checklist}

\begin{lstlisting}
□ Requirement Analysis
  □ Understand user need
  □ Identify AC
  □ Define dependencies

□ Design
  □ System design
  □ Database design
  □ UI mockups

□ Development
  □ Code implementation
  □ Code review
  □ Refactoring

□ Testing
  □ Unit testing
  □ Integration testing
  □ UAT

□ Documentation
  □ Code comments
  □ API docs
  □ User guide

□ Deployment
  □ Staging test
  □ Production deploy
  □ Monitor
\end{lstlisting}

\chapter{KẾT LUẬN}

Hệ thống quản lý khách sạn là một dự án phức tạp đòi hỏi:

\begin{itemize}
    \item \textbf{Phân tích kỹ lưỡng} các yêu cầu chức năng
    \item \textbf{Viết User Story rõ ràng} với AC cụ thể
    \item \textbf{Tổ chức Task hợp lý} theo Epic → Feature → US → Task
    \item \textbf{Ước tính Story Points} chính xác
    \item \textbf{Thực hiện kiểm thử toàn diện} ở mỗi cấp độ
    \item \textbf{Quản lý rủi ro} trong suốt dự án
\end{itemize}

Bằng cách tuân theo các nguyên tắc này, team có thể phát triển hệ thống chất lượng cao, đáp ứng đầy đủ yêu cầu, và được giao kịp thời.

\end{document}
